\documentclass[a4paper]{article}
\usepackage{graphicx} % Required for inserting images
\usepackage[font=it, width=0.9\linewidth]{caption}
\usepackage{mathtools}
\usepackage{listings}
\usepackage{graphicx}
\usepackage{bm}
\usepackage{caption}
\usepackage{float}
\usepackage{svg}
\renewcommand{\figurename}{Fig.}
\renewcommand{\tablename}{Tab.}

\title{Misura della caratteristica di due diodi a giunzione p-n}
\author{Bellini Samuele, Caprara Francesco}
\date{Novembre 2025, turno 3, tavolo 10}

\begin{document}

\maketitle

\section{Abstract}
\section{Introduzione}
\section{Apparato sperimentale e svolgimento}
Gli strumenti di misura che abbiamo utilizzato sono un oscilloscopio analogico GW~GOS-652 50Mhz e un multimetro 
digitale Iso-Tech~IDM~105. L'oscilloscopio ha una portata di 40~V, divisa in 10 sezioni definite dal valore scelto per
i volt associati ad ognuna delle 8 divisioni dello schermo.
I valori disponibili da associare alle divisioni sono:
\begin{table}[H]
    \centering
    % \setlength{\tabcolsep}{10pt}  %column separation%
    % \renewcommand{\arraystretch}{1.25} %row separation%
    \begin{tabular}{c  c  c  c  c  c  c  c  c  c  c }
        \textbf{Volts/Div} & 0.005 & 0.01 & 0.02 & 0.05 & 0.1 & 0.2 & 0.5 & 1 & 2 & 5 \\
    \end{tabular}
\end{table}
\noindent L'errore associato alle misure fatte con l'oscilloscopio è dato dalla formula
\[\sigma = \sqrt{{\sigma_L} ^2 + {\sigma_Z}^2 + {\sigma_C}^2 } \]
dove \(\sigma_L\) è l'errore sulla lettura, pari a \(\frac{1}{2} \frac{Volts/Div}{5}\), in quanto ogni divisione è
composta da 5 tacchette ed i segnali erano sufficientemente puliti da poter apprezzare la mezza tacchetta;
\(\sigma_Z\) è l'errore sullo zero, cioè l'errore dovuto all'imprecisione nel posizionare la linea di 0~V, ed è
pari alla \(\sigma_L\) associata agli 0.005~Volts/Div*****KAPE QUA DOBBIAMO DECIDERE SE DIRE DEI 20***, in quanto abbiamo utilizzato 
quella scala per posizionare lo zero, e \(\sigma_C\) è l'errore dichiarato dal costruttore, pari al 3\% della misura. 
\\Il multimetro è stato usato sia come voltmetro che come amperometro. Nella lettura di tensioni continue
ha una portata di 2000~V, nel range [-1000, 1000] (*****CREDO????*****). Nell'intervallo [0, 0.4]V l'errore da associare
alla misura è dato da 
\[\sigma = { 0.3\% \, misura + 0.2 mV } \]
dove \(misura\) è il valore letto dal multimetro in mV.
Nell'intervallo [0.4, 4]V, invece, l'errore da associare
alla misura è dato da 
\[\sigma = { 0.1\% \, misura + 2 mV } \]
Nella lettura di correnti continue la portata dell'amperometro è di 20~A, nel range [-10, 10]A. Nell'intervallo [0, 4]mA l'errore da associare
alla misura è dato da 
\[\sigma = { 0.4\% \, misura + 0.002 mA } \]
dove \(misura\) è il valore letto dal multimetro in mA.
Nell'intervallo [4, 400]mA, invece, l'errore da associare
alla misura è dato da 
\[\sigma = { 0.4\% \, misura + 0.02 mA }. \]
Abbiamo inoltre utilizzato il generatore da banco TTi EB2025T per ottenere una tensione costante di 5~V. Per poter fornire al circuito
un set di tensioni con step nell'ordine del mV abbiamo utilizzato un potenziometro da 1~k\(\Omega\) come partitore di tensione variabile, come verrà illustrato in seguito,
in quanto la risoluzione del generatore non sarebbe stata sufficiente.

Per verificare che oscilloscopio e multimetro fossero calibrati abbiamo realizzato il circuito riportato in \figurename~\ref{App_sper:circuiti} a sinistra
misurando sia con il multimetro che con l'oscilloscopio lo stesso punto del circuito, con una tensione che abbiamo fatto variare nel range [5, 800]mV tramite il potenziometro, che funge 
da resistenza variabile. Lo step di tensione tra i punti è pari al valore di una divisione dello schermo dell'oscilloscopio nel range in cui sono stati misurati, in modo da semplificare la lettura delle misure. 
Sono stati presi un totale di 37 punti.

Per misurare la curva di caratteristica in polarizzazione diretta di un diodo di silicio e uno di germanio abbiamo realizzato il circuito riportato in \figurename~\ref{App_sper:circuiti} a destra,
dove il multimetro funge da amperometro, misurando la corrente che passa attraverso il diodo, mentre l'oscilloscopio misura la tensione ai capi del diodo.
La tensione, per il germanio, è stata fatta variare, tramite il potenziometro, nel range [100, 760]mV, prendendo 41 coppie tensione-corrente, mentre per il germanio l'intervallo scelto è 
[20, 340]mV misurandone 65 coppie tensione-corrente.
\begin{figure}[H]
    \centering
    \includegraphics[width=.8\linewidth]{circuiti.png}
    \caption{A sinistra il circuito realizzato per calibrazione degli strumenti, a 
    destra il circuito realizzato per la misura della curva caratteristica dei diodi.\label{App_sper:circuiti}}
\end{figure}    
% \begin{table}[H]
%     \centering
%     \setlength{\tabcolsep}{10pt}  %column separation%
%     \renewcommand{\arraystretch}{1.25} %row separation%
%     \begin{tabular}{c | c | c }
%         \textbf{Volts/Div} & \textbf{Portata (V)} & \bm{{\(\sigma_L (mV)\)}} \\
%         \hline
%         0.005 & 0.04 &0.0005\\
%         0.01 & 0.08 & 0.001 \\
%         0.02 & 0.16 & 0.002 \\
%         0.05 & 0.4 & 0.005 \\
%         0.1 & 0.8 & 0.01  \\
%         0.2 & 1.6 & 0.02 \\
%         0.5 & 4 & 0.05 \\
%         1 & 8 & 0.1 \\
%         2 & 16 & 0.2 \\
%         5 & 40 & 0.5 \\
%     \end{tabular}
%     \caption{ Specifiche dell'oscilloscopio utilizzato
%     \label{tab:App_Sperim:dati_oscill}}
% \end{table}

\section{Risultati e discussione}
\section{Conclusioni}
\section{Bibliografia}
\section{Appendice}


\end{document}
